The module system is part of a larger compiler pipeline which implements the frontend.
An overview of this pipeline is given in \cref{fig:intro:structure}.

\begin{figure}[h]
  \begin{center}
    \tikzstyle{textonly}=[align=center, font=\sffamily]
    \tikzstyle{box}=[textonly, rectangle, draw=black, minimum width = 2.0cm]
    \tikzstyle{implemented}=[-{Classical TikZ Rightarrow[length=1.5mm]}, align=center, font=\sffamily]
    \tikzstyle{label}=[text=gray, align=center, font=\small\sffamily]

    \tikzmath{\distHoriz=2.0;}
    \tikzmath{\distHorizShort=1.8;}
    \tikzmath{\distHorizBot=3.0;}
    \tikzmath{\distVert=1.5;}
    \tikzmath{\distVertHalf=0.75;}

    \begin{tikzpicture}[scale=0.8, every node/.style={transform shape}]
      % Top Row
      \node (inputs) [box]                                        {List Unicode\\\small{\color{gray} Text Input}};
      \node (lexinput) [box, right = \distHoriz of inputs]            {List (Unicode $\times\ \mathbb{N}$)\\\small{\color{gray}Input to Lexer}};
      \node (tokens) [box, right = \distHorizShort of lexinput]           {List Token\\\small{\color{gray}Section \ref{subsec:lexer:tokens}}};
      \node (cst) [box, right = \distHorizShort of tokens]          {CST\\\small{\color{gray}Section \ref{subsec:parser:cst}}};
      % Bottom Row
      \node (rst) [box, below = \distVert of inputs]           {RST\\\small{\color{gray}Section X}};
      \node (rstop) [box, right = \distHorizBot of rst] {RST$_{\text{op}}$\\\small{\color{gray} Section X}};
      \node (rstbind) [box, right = \distHorizBot of rstop] {RST$_{\text{bind}}$\\\small{\color{gray} Section X}};

      \node[shape=coordinate,below = \distVertHalf of cst] (cornerone) {};
      \node[shape=coordinate,above = \distVertHalf of rst] (cornertwo) {};
    
    
      % Columnizer
      \draw[implemented] (inputs) to node (tolexinput) [above]{Columnizer} node [below, label] {Section \ref{sec:columnizer}} (lexinput);
      % Lexer
      \draw[implemented] (lexinput) to node (totokens) [above]{Lexer} node [below, label] {Section \ref{sec:lexer}} (tokens);
      % Parser
      \draw[implemented] (tokens) to node (tocst) [above]{Parser} node [below, label] {Section \ref{sec:parser}} (cst);
      % Renamer
      \draw[implemented] (cst.south) -- (cornerone) -- (cornertwo) node[pos=0.5,above] {Name Resolution} node[pos=0.5, below,label] {Section \ref{sec:module-system}} -- (rst.north);
      % Fixity
      \draw[implemented] (rst) to node (torstop) [above]{Fixity Resolution} node [below, label] {Section \ref{sec:operator-fixity}} (rstop);
      % Binding Groups
      \draw[implemented] (rstop) to node (torstbind) [above]{Binding Groups} node [below, label] {Section \ref{sec:dependency-order}} (rstbind);
    \end{tikzpicture}
  \end{center}
  \caption{Structure of the frontend.}
  \Description{Structure of the frontend.}
  \label{fig:intro:structure}
\end{figure}

For this project we are only concerned with the name resolution pass which implements the semantics of the module system.
To give a rough idea about the other elements of this pipeline:
The columnizer specifies the meaning of tab-stops according to the Haskell report and annotates each character in the input with its indentation level.
The Lexer specifies regular grammars, the syntax of tokens and gives a formal specification of the maximum munch rule following \citet{Egolf2021verbatim}.
The parser parses indentation sensitive syntax following \citet{Adams2013parsing} and \citet{Adams2014parsec}.
Fixity resolution does the proper disambiguation of expressions involving operators, for example $a + b * c$ and takes custom infix declarations into account.
This has to happen after name resolution, due to how the Haskell report specifies it.
Lastly, mutually recursive bindings groups are computed for bindings and toplevel declarations. This is required for the correct kind and type inference.